%%%%%%%%%%%%%%%%%%%%%%%%%%%%%%%%%
\section{Introdução}

Este trabalho tem como objetivo apresentar o desenvolvimento do Arianagram, uma rede social minimalista baseada no modelo cliente-servidor, utilizando comunicação via sockets TCP e suporte a múltiplas conexões simultâneas. A proposta central do sistema é permitir que diferentes usuários se conectem a um servidor, publiquem mensagens, sigam outros usuários e recebam notificações em tempo real quando os perfis acompanhados realizarem novas postagens.

Ao longo do desenvolvimento, foram aplicados conceitos importantes de redes de computadores, como criação e configuração de sockets TCP, uso de IPv4 e IPv6, envio e recebimento de mensagens binárias de tamanho fixo, tratamento de conexões encerradas, sincronização de estruturas compartilhadas e comunicação assíncrona entre servidor e clientes.

\section{Arquitetura do Sistema}

\subsection{Visão Geral}

Nesse sistema, o servidor concentra a lógica principal da aplicação, enquanto os clientes funcionam como interface de interação, enviando comandos e exibindo mensagens recebidas.

A arquitetura do servidor segue o modelo Single-Process, Multi-Thread. O processo principal permanece aguardando novas conexões e, a cada cliente conectado, cria uma thread dedicada para atender aquela sessão. Isso permite que múltiplos usuários interajam simultaneamente com o sistema.

Além do fluxo tradicional de requisição e resposta, o Arianagram também implementa notificações push, explorando o paradigma Publish-Subscribe. Assim, quando um usuário publica uma mensagem, o servidor verifica seus seguidores ativos e envia a postagem diretamente para os clientes conectados que seguem aquele perfil.

\subsection{Estruturas de Dados}

No servidor, foram utilizadas estruturas de dados simples para armazenar as informações necessárias ao funcionamento do sistema. As principais estruturas estão relacionadas ao feed global de mensagens e ao controle de seguidores.

A estrutura do feed, implementada em \texttt{feed.c}, funciona como um buffer circular. Ela armazena as mensagens publicadas pelos usuários, mantendo um contador de mensagens válidas e uma posição inicial para indicar a mensagem mais antiga. Quando o limite de armazenamento é atingido, as mensagens mais antigas passam a ser sobrescritas pelas novas publicações. Essa abordagem permite manter o histórico recente de mensagens de forma simples e eficiente.

Já a estrutura de seguidores, implementada em \texttt{followers.c}, mantém a relação entre usuários seguidos e seus respectivos seguidores. Para cada usuário, é armazenada uma lista de perfis que o seguem. Assim, quando um usuário realiza uma postagem, o servidor consulta essa estrutura para identificar quais clientes devem receber uma notificação \texttt{MSG\_PUSH}.

Além disso, os nomes de usuário são normalizados antes do armazenamento, removendo o caractere @. Essa decisão facilita as comparações entre usuários e também permite tratar corretamente casos como seguidores duplicados e auto-follow, que são ignorados silenciosamente pelo sistema.

\subsection{Funcionamento da Thread no Servidor}

A arquitetura do servidor segue o modelo Single-Process, Multi-Thread. O processo principal permanece aguardando conexões e, sempre que um novo cliente é aceito com \texttt{accept()}, cria uma thread dedicada para atender aquela conexão.

\begin{figure}[H]
  \centering
  \includegraphics[width=0.8\textwidth]{thread_life_cycle.png}
  \caption{Ciclo de vida de uma thread no servidor}
  \label{fig:thread_life_cycle}
\end{figure}

O uso de \texttt{pthread\_detach()} permite que os recursos da thread sejam liberados automaticamente ao final da execução, sem a necessidade de utilizar \texttt{pthread\_join()} no processo principal.

Dentro da função \texttt{client\_thread}, a thread recebe inicialmente a mensagem \texttt{MSG\_CONNECT}, enviada pelo cliente logo após a conexão. A partir dessa mensagem, o servidor identifica o usuário conectado e registra esse cliente na lista de clientes ativos.

\begin{figure}[H]
  \centering
  \includegraphics[width=0.8\textwidth]{client_connected.png}
  \caption{Lista de clientes ativos no servidor}
  \label{fig:active_clients}
\end{figure}

Depois do handshake, a thread entra em um loop de processamento, permanecendo ativa enquanto a conexão com o cliente estiver aberta. Nesse loop, cada mensagem recebida é enviada para \texttt{handle\_message()}, que trata comandos como \texttt{MSG\_POST}, \texttt{MSG\_FOLLOW} e \texttt{MSG\_READ}.

\begin{figure}[H]
  \centering
  \includegraphics[width=0.8\textwidth]{message_handle.png}
  \caption{Processo de leitura de mensagens}
  \label{fig:leitura_mensagens}
\end{figure}

No caso de uma mensagem \texttt{MSG\_POST}, a thread armazena a publicação no feed global e aciona o envio de notificações push para os seguidores ativos daquele usuário.

\begin{figure}[H]
  \centering
  \includegraphics[width=0.8\textwidth]{post_store.png}
  \caption{Processo de leitura de armazenamento e push de postagens}
  \label{fig:armazenamento_push}
\end{figure}

Quando o cliente encerra a conexão, \texttt{recv\_all()} detecta o fechamento do socket, a thread sai do loop, remove o cliente da lista de ativos, fecha o socket e finaliza sua execução.

\begin{figure}[H]
  \centering
  \includegraphics[width=0.8\textwidth]{close_thread.png}
  \caption{Fechamento da thread}
  \label{fig:fechamento_thread}
\end{figure}

Assim, o ciclo de vida da thread pode ser resumido em: criação após \texttt{accept()}, identificação do cliente via \texttt{MSG\_CONNECT}, registro como cliente ativo, processamento dos comandos recebidos e limpeza dos recursos ao final da conexão.

\section{Desafios e Dificuldades}

O principal desafio desse sistema ocorreu principalmente no envio de pushes para permitir que uma thread enviasse mensagens para sockets associados a outros clientes. Quando um usuário realiza um \texttt{MSG\_POST}, a thread responsável por essa conexão precisa identificar quais seguidores estão ativos e encaminhar uma mensagem \texttt{MSG\_PUSH} diretamente para os sockets correspondentes.

Para isso, foi necessário manter uma lista de clientes ativos, relacionando cada usuário conectado ao seu socket. Além disso, como essa lista, o feed global e a tabela de seguidores são compartilhados entre múltiplas threads, foi necessário proteger o acesso a essas estruturas com mutexes, evitando inconsistências causadas por acessos concorrentes.

Outro ponto importante foi tratar corretamente desconexões. Quando um cliente encerra a conexão, ele deve ser removido da lista de ativos para impedir que o servidor tente enviar pushes para sockets inválidos.

Por fim, no lado do cliente, foi utilizado \texttt{select()} para monitorar simultaneamente o teclado e o socket, permitindo que notificações \texttt{MSG\_PUSH} fossem recebidas de forma assíncrona, mesmo quando o usuário não estivesse digitando comandos.

\section{Conclusão}

Assim, concluímos que este projeto permitiu a aplicação prática de conceitos de redes de computadores, programação em C e desenvolvimento de sistemas baseados no modelo cliente-servidor. A implementação do sistema \texttt{Arianagram} proporcionou um estudo mais aprofundado sobre comunicação via sockets TCP, múltiplas conexões simultâneas, uso de threads e sincronização de estruturas compartilhadas.

Além disso, o projeto possibilitou compreender melhor o funcionamento de um sistema baseado em Publish-Subscribe, no qual o servidor não apenas responde às requisições dos clientes, mas também envia notificações assíncronas por meio de mensagens \texttt{MSG\_PUSH}. Dessa forma, o desenvolvimento contribuiu para consolidar conhecimentos sobre concorrência, gerenciamento de conexões e comunicação em tempo real entre processos distribuídos.